\documentclass[a4paper, 11pt]{article}

% --- Packages ---
\usepackage[utf8]{inputenc}
\usepackage[T1]{fontenc}
% \usepackage[french]{babel}
\usepackage[margin=1.5cm, top=2cm, bottom=2cm]{geometry}
\usepackage{xcolor}
\usepackage{fancyhdr}
\usepackage{titlesec}
\usepackage{enumitem}
\usepackage{array}
\usepackage{longtable}
\usepackage{booktabs}
\usepackage{colortbl}
\usepackage{graphicx}

% --- Définition des couleurs ---
\definecolor{primaryBlue}{RGB}{20, 50, 90}
\definecolor{lightBlue}{RGB}{235, 240, 247}

% --- Styles de document ---
\pagestyle{fancy}
\fancyhf{}
\fancyfoot[L]{\footnotesize \textcolor{gray}{Projet Wakala - Étude Utilisateurs (11 répondants)}}
\fancyfoot[R]{\footnotesize \textcolor{gray}{Page \thepage}}
\renewcommand{\headrulewidth}{0pt}
\renewcommand{\footrulewidth}{0.4pt}

% --- Personnalisation des titres ---
\titleformat{\section}{\Large\bfseries\color{primaryBlue}}{}{0em}{}[\titlerule]
\titleformat{\subsection}{\large\bfseries\color{primaryBlue}}{}{0em}{}

\newcolumntype{L}[1]{>{\raggedright\arraybackslash}p{#1}}

\begin{document}

% --- En-tête ---
\begin{center}
    {\Huge \textcolor{primaryBlue}{\textbf{WAKALA}}}\\[0.5cm]
    {\LARGE \bfseries \textcolor{primaryBlue}{Synthèse de l'Étude Utilisateurs}}\\[0.2cm]
    {\large \textit{Analyse des attentes (UX) et des craintes lors de l'achat de voitures d'occasion}}
\end{center}

\vspace{0.8cm}

% --- Introduction ---
\section*{Introduction}
Ce document présente d'abord la simulation détaillée des réponses au questionnaire administré à 11 participants, suivie d'une synthèse stratégique concernant leurs attentes et craintes lors de l'achat de véhicules d'occasion.

\vspace{0.5cm}

% --- 1. Simulation des Réponses Détaillées ---
\section*{1. Simulation des Réponses (Les 11 Profils)}

\subsection*{Profils et Habitudes d'achat (Questions 1 à 5)}
\begin{footnotesize}
\begin{longtable}{L{0.5cm} L{1.8cm} L{2.5cm} L{2cm} L{2.2cm} L{2.2cm} L{2.5cm}}
\toprule
\rowcolor{primaryBlue}
\textcolor{white}{\textbf{N°}} & \textcolor{white}{\textbf{Âge}} & \textcolor{white}{\textbf{Profession}} & \textcolor{white}{\textbf{Ville}} & \textcolor{white}{\textbf{Déjà acheté ?}} & \textcolor{white}{\textbf{Où ?}} & \textcolor{white}{\textbf{Durée rech.}} \\
\midrule
\endhead
1 & 22 ans & Étudiant & Casablanca & Non & N/A & 1 mois \\
\rowcolor{lightBlue}
2 & 24 ans & Étudiante & Rabat & Oui & Avito & 3 mois \\
3 & 45 ans & Employé & Meknès & Oui & Avito / Garage & 2 mois \\
\rowcolor{lightBlue}
4 & 52 ans & Professeur & Rabat & Oui & Concession & 1 mois \\
5 & 58 ans & Retraité & Casablanca & Oui & Entre amis & 4 mois \\
\rowcolor{lightBlue}
6 & 25 ans & Employé & Casablanca & Oui & Avito & 2 mois \\
7 & 48 ans & Femme au foyer & Meknès & Non & N/A & 2 semaines \\
\rowcolor{lightBlue}
8 & 21 ans & Étudiant & Rabat & Non & N/A & 2 mois \\
9 & 55 ans & Retraité & Meknès & Oui & Garage local & 1 mois \\
\rowcolor{lightBlue}
10 & 49 ans & Professeur & Casablanca & Oui & Avito & 3 mois \\
11 & 23 ans & Employée & Rabat & Oui & Avito & 2 mois \\
\bottomrule
\end{longtable}
\end{footnotesize}

\subsection*{Expérience Utilisateur (UX) et Craintes (Questions 6 à 9)}
\begin{itemize}[label=\textbullet, leftmargin=*]
    \item \textbf{Q6. Qu'est-ce qui fait dire "cette annonce est bien faite" ?} \\
    \textit{Réponses :} "Beaucoup de photos nettes intérieur/extérieur" (7/11), "Description honnête sans fautes" (3/11), "Prix final clair" (1/11).
    
    \item \textbf{Q7. Quelles informations sont indispensables ?} \\
    \textit{Réponses :} Historique des accidents (11/11), Kilométrage certifié (11/11), État du moteur (9/11), Nombre de propriétaires (6/11).
    
    \item \textbf{Q8. Préfères-tu chercher seul(e) ou être guidé(e) ?} \\
    \textit{Réponses :} Seul(e) sur internet pour fouiller (8/11), mais guidé par un expert physique pour la visite finale (10/11). L'aide par chatbot intéresse surtout les 21-25 ans.
    
    \item \textbf{Q9. Quelle est ta plus grande difficulté lors de la recherche ?} \\
    \textit{Réponses :} Ne pas savoir quel modèle est adapté à ma vie (9/11), Ne pas comprendre le jargon technique (7/11), Se perdre parmi trop d'annonces (4/11).
\end{itemize}

\vspace{0.8cm}

% --- 2. Synthèse Globale ---
\newpage
\section*{2. Synthèse Globale : Ce qu'il faut retenir}

Les réponses de notre panel montrent une chose très claire sur le marché marocain : les acheteurs veulent utiliser internet pour trouver une voiture, mais ils sont complètement perdus face à la complexité des choix. En effet, si la majorité des personnes (8 sur 11) aiment faire leurs recherches seules sur internet pour comparer les offres, le moment de sélectionner le modèle exact parmi des centaines génère un fort stress (peur de se tromper). Presque tout le monde veut alors être accompagné ou conseillé de manière personnalisée, mais les sites actuels abandonnent l'utilisateur exactement à ce moment-là. Cette anxiété s'explique par la difficulté d'estimer si un modèle correspond vraiment à son style de vie (cité 9 fois) et à son budget réel. 

Aujourd'hui, de simples filtres de recherche ne suffisent plus. C'est pour répondre à ce besoin précis que Wakala intègre un \textbf{Chatbot IA} pour conseiller l'acheteur comme un véritable ami, et un \textbf{Moteur de Recommandation} intelligent pour lui suggérer la voiture parfaite. 
Si un utilisateur a peur des arnaques au kilométrage, le Chatbot s'adapte, le rassure, et lui explique comment vérifier les papiers. S'il hésite entre deux modèles, le Chatbot résume les points forts de chacun de manière très simple (ex: "La Clio coûtera moins cher en assurance, la Peugeot a un plus grand coffre"). L'Intelligence Artificielle est ici utilisée pour remettre de l'humain et de l'empathie dans un processus d'achat froid et technique.



\end{document}
